\documentclass[10pt,a4paper]{amsart}
\usepackage[margin=19mm]{geometry}
\usepackage{amsmath,amssymb,mathtools}
\usepackage[scr=rsfs,cal=euler]{mathalfa}
\usepackage{xfrac}
\usepackage[unicode,hidelinks,bookmarks=false]{hyperref}
\newcommand{\ie}{i.e.\ }
\newcommand{\cf}{cf.\ }
\newcommand{\resp}{respectively\ }
\newcommand{\Subsection}{\subsection}
\providecommand{\nobreakdash}{\nobreak\hbox{-}}
\setlength{\emergencystretch}{2em}
\multlinegap0pt
\usepackage{amssymb}
\usepackage{enumerate}
\usepackage{mathtools}
\usepackage{xcolor}
\newtheorem{lemma}{Lemma}
\title{\textbf{Characterizations of Almost Ricci Bourguignon Solitons}}
\author{Mohammad Aqib, Hemangi Madhusudan Shah, Dhriti Sundar Patra}
\date{Source: arXiv:2512.17704v1 (CC BY 4.0)}
\begin{document}
\maketitle

\noindent\emph{Source and licence.} \textbf{Characterizations of Almost Ricci Bourguignon Solitons}. Version arXiv:2512.17704v1 (CC BY 4.0), distributed by arXiv under the Creative Commons Attribution 4.0 International licence, http://creativecommons.org/licenses/by/4.0/, as recorded in the arXiv licence metadata for this exact version and shown on the abstract page https://arxiv.org/abs/2512.17704v1. Full licence notice: \texttt{licenses/arxiv-2512.17704-CC-BY-4.0.txt}.\par\smallskip

\noindent\textit{Editorial scope.} Counted unit: the article's lemma lem2 (source0.tex lines 473--482). The article's complete original proof of it is delivered with it (source0.tex lines 484--511, one \texttt{proof} environment of 28 lines). No mathematics is written, repaired or re-derived: the statement and the proof are the article's own lines, in the article's own notation, and no retained line is substituted or re-flowed.  Retained uncounted as the source-local context the proof needs: lines 412--414; lines 445--450.  Nothing was omitted from any retained span, so the package carries no omission record.  No retained line contains a citation, so the wrapper carries no bibliography and no bibliography entry.  No retained line contains a figure environment, an image-inclusion command or drawing code, so no image is delivered and the package declares no asset dependency.  Editorial formatting only, in addition: an AMS \texttt{amsart} preamble that restates the article's own declarations and macros used by the retained lines, namely the environments \texttt{lemma} on separate counters, together with the standard packages those lines need.  The retained environments print in this document as Lemma 1 in order of appearance, whereas the article prints the same items with the numbering of its own full document.  The complete native source of the article as distributed in the arXiv e-print for this exact version is bundled unchanged as \texttt{sources/191338/source0.tex}.
\begin{equation}\label{div}
    \operatorname{div} \xi=(n (\lambda +\rho S)-S).
\end{equation}

\begin{lemma}\label{lem1}
    Let $(M^n, g, \xi, \lambda,\rho)$ be a compact almost RB-soliton. Then
\begin{equation}\label{e1}
\int_{M}\left((\lambda+\rho S)S -\|Q\|^{2}-\frac{1}{2} S(n (\lambda +\rho S)-S)\right)=0.
\end{equation}
\end{lemma}

\begin{lemma}\label{lem2}
     Let $(M^n, g, \xi, \lambda,\rho)$ be a compact almost RB-soliton. Then
% \begin{equation}\label{10.2}
%     \int_M\left|\operatorname{Ric}-\frac{S}{n} g\right|^{2} =
%     \frac{(n-2)}{2n(1-n\rho)}\int_M ng(\nabla\lambda,\xi)-g(\nabla \operatorname{div}(\xi),\xi)\quad \mbox{for } \rho \neq \frac{1}{n}.
%     \end{equation}
    \begin{equation}\label{10.3}
        \int_M\left|\operatorname{Ric}-\frac{S}{n} g\right|^{2} =  \frac{(n-2)}{2n}\int_M g(\nabla S,\xi).
    \end{equation}
\end{lemma}

\begin{proof}

    We have
    \begin{equation}\label{ricsbyn}
        \left|\operatorname{Ric}-\frac{S}{n} g\right|^{2}=\|Q\|^{2}-\frac{S^{2}}{n}.
    \end{equation}
Now \eqref{div} yields that,
$$
\frac{1}{n} \operatorname{div}(S \xi)=\frac{1}{n} \xi(S)+ S(\lambda +\rho S)-\frac{S^{2}}{n}.
$$
Thus,
$$
\left|R i c-\frac{S}{n} g\right|^{2}=\|Q\|^{2}+\frac{1}{n} \operatorname{div}(S \xi)-\frac{1}{n} \xi(S)-S(\lambda + \rho S).
$$
This implies that,
$$
\int_{M}\left|R i c-\frac{S}{n} g\right|^{2}=\int_{M}\left(\|Q\|^2-\frac{1}{n}\xi(S)-S(\lambda+\rho S)\right) .
$$
Applying Lemma \ref{lem1}, we conclude
$$
\int_{M}\left|R i c-\frac{S}{n} g\right|^{2}=\int_{M}\left(\frac{S^2}{2}-\frac{1}{n}\xi(S)-\frac{1}{2}S(n(\lambda+\rho S)\right) .
$$
Now using  $\frac{1}{2} \operatorname{div}(S \xi) = \frac{1}{2}\xi(S)+ \frac{S}{2} \operatorname{div} \xi$ and  \eqref{div} in the above equation, we infer
    \begin{equation}\label{10}
    \int_M\left|\operatorname{Ric}-\frac{S}{n} g\right|^{2}=\frac{(n-2)}{2n}\int_M g(\nabla S,\xi).
\end{equation}
% Now, if we use the almost RB-soliton equation on the right-hand side of the above equation, we obtain \eqref{10.2} as stated.
\end{proof}

\end{document}
